\begin{abstract}
We prove that every nonempty finite oriented graph has a vertex with at least
as many vertices at directed distance two as at directed distance one.
This resolves Seymour's second neighborhood conjecture positively.
\end{abstract}
\section{Introduction}\label{sec:introduction}

An \emph{oriented graph} is a finite directed graph with no loops, no multiple
arcs, and no pair of opposite arcs.  Its underlying undirected graph need
not be complete or connected.  For a vertex \(v\), write
\[
 N_1^+(v)=\{u:v\to u\},\qquad
 N_2^+(v)=\{w\notin\{v\}\cup N_1^+(v):
                 \text{ there is a vertex }u\text{ with }v\to u\to w\}.
\]
Thus \(N_2^+(v)\) consists of the vertices at directed distance exactly two
from \(v\).  Each vertex is counted once, regardless of the number of paths
reaching it.  Seymour's second neighborhood conjecture asserts that some
vertex satisfies \(\abs{N_1^+(v)}\leq\abs{N_2^+(v)}\).

\begin{theorem}\label{thm:main}
Every nonempty finite oriented graph has a vertex \(v\) such that
\[
 \abs{N_1^+(v)}\leq \abs{N_2^+(v)}.
\]
\end{theorem}

Theorem~\ref{thm:main} resolves the conjecture positively, without a
connectivity or degree hypothesis.  In particular, a sink already satisfies
the conclusion.  The conjecture is a basic comparison between the first two
out-neighborhoods: removing vertices already reached in one step is essential
to its assertion.

The theorem also gives the established vertex- and arc-weighted forms of
the conjecture through Seacrest's equivalences~\cite{Seacrest2015}.
The vertex-weighted conclusions have two different quantifier orders:
for every nonnegative weighting there is a suitable vertex, and there is
a weighting of total mass one for which every vertex is suitable.
Section~\ref{sec:weighted-consequences} states these forms precisely.
Section~\ref{sec:directed-cycle-consequences} records consequences for
